\subsection{The quaternion cover: pointed reduction and the proper fixed group}

\begin{statement}{A.35}{Proved}{The quaternion cover has no proper fixed image at odd primes}
Let \(C=C_Q\) be the smooth quartic \(X^4+Y^4=2Z^4\), and
let \(S=\{Z=0\}\). After a finite extension the pair \((C,S)\)
has the cyclic description
\[
             U^8=-64B(B-1)^2,\qquad
             C\setminus S\longrightarrow
                    \mathbf P^1\setminus\{0,1,\infty\}.
                                                        \tag{A.35.1}
\]
This is an \'etale cyclic cover of degree eight. There are respectively
one, two, and one points above \(0,1,\infty\), with ramification
indices \(8,4,8\).

In the unpointed stable reduction at \(2\), the three elliptic
components meet a central rational component. All four points of
\(S\) specialize to the same node: the node adjacent to the unique
elliptic component on which the central involution of the original
\(Q_8\)-cover acts trivially. In particular, the four marked points
are not distinct smooth points of that central component.

Let \(P\in S\), let \(k\) be a number field over which \(P\) and
the curve are defined, and let \(\bar s_P\) be the section of the
proper arithmetic fundamental group defined by \(P\). For every
odd prime \(\ell\) and every open subgroup
\(\Delta_0\subset\operatorname{Gal}(\bar k/k^{\mathrm{ab}})\),
\[
 \left(\pi_1(C_{\bar k},P)^{(\ell)}\right)^{
                                  \bar s_P(\Delta_0)}=1.
                                                        \tag{A.35.2}
\]
Consequently every fixed element in the corresponding punctured
pro-\(\ell\) group, for every tangent packet at \(P\), belongs to
the closed normal subgroup generated by the four cusp inertias.
This consequence does not identify that punctured fixed subgroup
with the selected inertia; the normal subgroup in question is
much larger than a single inertia group.
\end{statement}

\begin{proof}
\paragraph{The cyclic model and the specialization of the cusps.}
Over a finite extension, put the quartic in the form
\(W^4=1-T^4\), with \(S=\{W=0\}\). Set
\[
 z=\frac{T-1}{i(T+1)},\qquad B=z^2,\qquad
 U=W(1-iz).
\]
Then \(T=(1+iz)/(1-iz)\) and
\[
 (1-T^4)(1-iz)^4=8iz(z^2-1).
\]
Squaring gives (A.35.1). Conversely \(z=U^4/(8i(B-1))\),
so these formulas identify the function fields. The Kummer
exponents at \(0,1,\infty\) are \(1,2,-3\), giving the stated
indices and numbers of points. The four ramification points are
exactly \(W=0\). The program
\texttt{code/fgh\_q8\_models.py} checks the displayed identities
over \(\mathbf Q(i)\).

The unpointed stable fibre with a central rational component
and three elliptic tails is the explicit computation of
\cite[Example 10.4.39, pp.\ 548--549]{fghLiu};
it is also recorded in
\cite[Table 7.1, p.\ 50]{fghQuarticReduction}.
Only that established Fermat computation is used here, not
the latter paper's conjectural general reduction criteria.
The elliptic tails are supersingular: the good special fibre
of the Jacobian is their product and is isogenous to a power
of the supersingular reduction of the CM elliptic curve
with \(j=1728\).

Let \(\tau\) generate the cyclic group in (A.35.1), and put
\(\sigma=\tau^4\). In the quartic coordinates
\(\tau^2\) acts by \(W\mapsto iW\), while
\(\sigma\) acts by \(W\mapsto-W\).
Thus \(C/\langle\tau^2\rangle\) has genus zero and
\(C/\langle\sigma\rangle\) has genus one. On rational
prime-to-two Tate homology, the \(\tau^2\)-invariants are
therefore zero and the \(\sigma\)-invariants have dimension two.

The permutation of the three tails induced by \(\tau\) is a
transposition. Indeed its order divides eight, so the only
alternatives are a transposition and the identity. In the latter
case it would act on each elliptic tail as an automorphism
fixing its attachment point. Such an automorphism cannot have
order eight: a finite-order elliptic endomorphism has an
integral quadratic characteristic polynomial, whereas a
primitive eighth root has degree four over \(\mathbf Q\).
Its fourth power would therefore be the identity on each
tail. This contradicts the dimension of the
\(\sigma\)-invariant homology just calculated.

Write \(E_0\) for the fixed tail and \(E_1,E_2\) for the
exchanged tails. On \(E_0\), the action of \(\tau\) has order
four: its fourth power is the identity, and its square cannot
have invariant homology. Its only geometric fixed point is
the attachment point. To see the uniqueness, take that point
as origin. The endomorphism \(\tau-1\) has degree two, and a
supersingular elliptic curve in characteristic two has no
nonzero geometric two-primary torsion. On the central
rational component, \(\tau\) exchanges two of the three
attachment points and fixes the third. A nonidentity
involution of \(\mathbf P^1\) in characteristic two has a
unique fixed point, namely this third attachment point.

The cusps above \(B=0,\infty\) are fixed by \(\tau\). Their
specializations cannot lie on the two exchanged tails, and
the preceding fixed-point calculations force them to the
node adjacent to \(E_0\). Finally the original \(Q_8\) deck
group acts transitively on \(S\) and commutes with its central
involution \(\sigma\). This involution acts trivially on
\(E_0\) and as minus one on each of the other two tails:
the fixed-tail contribution already exhausts its
two-dimensional invariant homology. The \(Q_8\) action
therefore preserves \(E_0\) and its attachment node.
Transitivity puts all four cusp specializations at that node.
This proves the claimed pointed correction without requiring
any assertion about the full graph permutation of a coordinate
sign involution.

\paragraph{The actual based monodromy.}
Retain just the one marked point \(P\). In the stable pointed
model, its specialization lies on a rational bridge inserted
between \(E_0\) and the original central component. Let the
two new edge lengths be \(a,b>0\), with \(a\) on the \(E_0\)
side, and let \(d_1,d_2>0\) be the two remaining elliptic-arm
lengths. These inequalities are strict: locally the original
node has equation \(uv=\varpi^e\), and a section through
the node has \(v(u(P)),v(v(P))>0\); these are precisely the
two lengths after a finite extension. Rational components
of valence two may be contracted, adding their lengths.

Put \(V_i=H_1(E_i,\mathbf Q_\ell)\), choose symplectic bases
\(x_i,y_i\), and write \(\omega_i=[x_i,y_i]\).
The rational graded Lie algebra of the proper curve is
\[
 \mathfrak s=
 \mathbb L(V_0\oplus V_1\oplus V_2)/
                  (\omega_0+\omega_1+\omega_2).
                                                        \tag{A.35.3}
\]
The based nonabelian Picard--Lefschetz formula
\cite[Definition 3.1.5 and Theorem 3.1.6]{fghBettsDogra},
applied with the base point on the displayed bridge, gives
the first monodromy derivation
\[
 \begin{split}
 N(v_0)&=a[\omega_0,v_0],\\
 N(v_i)&=[b(\omega_1+\omega_2)+d_i\omega_i,v_i],
                                      \qquad i=1,2.
 \end{split}                                             \tag{A.35.4}
\]
The term with coefficient \(b\) is the twist around the side
containing both \(E_1,E_2\). Thus (A.35.4) is a based formula;
an unspecified inner term is not being suppressed.
Equivalently, set \(m_0=a+b\), \(m_i=d_i\) for \(i=1,2\),
and \(w=-b\omega_0\). Then
\[
       N(v_i)=[w+m_i\omega_i,v_i].
                                                        \tag{A.35.5}
\]
All nodes are separating, so the degree-one monodromy
derivation vanishes and \(N\) has degree two. Filling the
chosen cusp kills every change of tangential packet, leaving
the proper point section used in this formula.

\paragraph{An unbounded kernel calculation in the surface quotient.}
We prove, in the actual quotient (A.35.3), that
\[
                    \ker N=\mathbf Q_\ell\omega_0.
                                                        \tag{A.35.6}
\]
Here is a normal form that makes the quotient issue explicit.
Present its enveloping algebra using \(u_1,u_2,x_i,y_i\),
put \(u_0=-u_1-u_2\), and impose the three relations
\[
                    x_i y_i=y_i x_i+u_i,\quad i=0,1,2.
                                                        \tag{A.35.7}
\]
The leading words \(x_i y_i\) have no self-overlaps or mutual
overlaps. Ordering the \(x_i\)'s above the \(y_i\)'s and the
\(u_j\)'s gives terminating reduction (give \(u_j\) degree two).
Consequently the words avoiding all three leading words
form a basis. In particular
\(B=\mathbf Q_\ell\langle u_1,u_2\rangle\) embeds.

Group these normal words by the ordered sequence of their
tail letters \(x_i,y_i\). Between consecutive tail letters
are slots containing words in \(B\). If consecutive tail
letters form \(x_i y_i\), their separating slot must be
nonempty; the other slots are unrestricted. The derivation
inserts only \(u\)-letters, so it preserves every such
normal-word component. On \(B\) it is \(\operatorname{ad}w\).
In a component with at least one tail letter, let its first
tail letter have index \(i\). On the first \(B\)-slot the
operator that increases that slot's degree is exactly
\[
                      L_w+m_iR_{u_i}.
                                                        \tag{A.35.8}
\]
For instance, \(N(f)=[w,f]\) and the left part of
\(N(x_i)=[w+m_i u_i,x_i]\) together give
\(wf+fm_i u_i\). Contributions from later slots do not
increase the first-slot degree.

The operator (A.35.8) is injective. For \(i=0\) it is
\(-bL_{u_0}+(a+b)R_{u_0}\). Choose \(u_0\) and a transverse
linear letter as free generators of \(B\). On pure powers
of \(u_0\), the coefficient is \(a\ne0\). Among words
containing the transverse letter, choose the greatest
initial \(u_0\)-prefix length: left multiplication increases
it, right multiplication does not, and \(-b\ne0\).
For \(i=1,2\), the linear letters \(w,u_i\) are independent.
Count occurrences of \(w\); the left summand increases
their number, while the right summand does not.
Taking the greatest first-slot degree in any finite
linear combination now proves that a component with a
tail letter has zero kernel.

Thus the associative kernel lies in \(B\), where it is
the centralizer of the linear letter \(w\), namely
\(\mathbf Q_\ell[w]\). The elements \(u_1,u_2\) are
primitive, so \(B\) is an embedded Hopf algebra. In
characteristic zero the only primitive polynomials in
one primitive letter are its scalar multiples. Intersecting
with \(\mathfrak s\) proves (A.35.6). This argument does not
pass a free-algebra kernel through the surface relation.
The certificate \texttt{code/a\_Q8\_proper\_fixed.py}
checks the absence of overlap ambiguities and the slot
matrices through degree nine for positive sample lengths;
the degree-independent proof is (A.35.8).

\paragraph{A separated integral filtration.}
For completeness, we use an integral augmentation filtration,
so no rational Lie element is silently lifted to a group
element. Write \(P_\ell=\pi_1(C_{\bar k},P)^{(\ell)}\).
Its pro-\(\ell\) presentation has six generators and the
single relation \(\prod_i[\alpha_i,\beta_i]=1\).
In \(\mathbf Z_\ell\langle\!\langle X_i,Y_i\rangle\!\rangle\),
impose the exact relation
\[
        \prod_i[1+X_i,1+Y_i]-1=0.
                                                        \tag{A.35.9}
\]
Its initial quadratic part is \(\sum_i[X_i,Y_i]\), up to
one harmless sign convention. Orient the coefficient-one
term \(X_0Y_0\) as a reduction rule. This word has no
self-overlap. At each finite degree cutoff, reduction
terminates: higher-degree terms decrease the truncated
order, and at a fixed degree the remaining quadratic
terms are lexicographically smaller. Disjoint reductions
commute; there is no overlap ambiguity. The compatible
normal-word bases therefore give a complete quotient
algebra \(\mathcal A\), free over \(\mathbf Z_\ell\) in
each graded degree, with
\[
 \operatorname{gr}\mathcal A=
 \mathbf Z_\ell\langle X_i,Y_i\rangle/
                       \left(\sum_i[X_i,Y_i]\right).
                                                        \tag{A.35.10}
\]
This is also the primitive one-relator situation of
\cite[Theorem 2 and its corollary, pp.\ 142--143]{fghLabute1967};
the direct monic reduction establishes what is needed here
without restricting the prime.

The map \(P_\ell\to\mathcal A^\times\) is injective.
For this continuous group map, give the power-series algebra
its compact \((\ell,X_i,Y_i)\)-adic topology. Its relevant
unit subgroup is an inverse limit of finite \(\ell\)-groups.
This compact topology is distinct from the degree filtration
used to extract first terms.
Indeed every finite \(\ell\)-group quotient \(Q\) permits
evaluation in \(\mathbf Z_\ell[Q]\). Its augmentation ideal
is topologically nilpotent for the \(\ell\)-adic topology,
since the augmentation ideal in \(\mathbf F_\ell[Q]\) is
nilpotent: some power is contained in
\(\ell\mathbf Z_\ell[Q]\), and iterating gives convergence.
Hence all the displayed series converge under
evaluation, and a nonidentity image in \(Q\) is detected
in its group algebra. Finite \(\ell\)-quotients separate
the points of \(P_\ell\).
The induced augmentation filtration on \(P_\ell\) is
therefore separated. Every first nonzero class remains
nonzero after tensoring with \(\mathbf Q_\ell\), and is
primitive for the associated graded Hopf coproduct.
Indeed the closed relation ideal is a Hopf ideal, since
the word set equal to one in (A.35.9) is group-like.
The primitive subspace of
\(\operatorname{gr}\mathcal A\otimes\mathbf Q_\ell
 =U(\mathfrak s)\) is exactly \(\mathfrak s\).
Thus (A.35.6) applies to the first class of every actual
nonidentity pro-\(\ell\) element.

\paragraph{From the local kernel to the global fixed group.}
Enlarge the number field to \(L\), with
\(G_L\cap\Delta_k\subset\Delta_0\), so that the stable model,
the components, and the point are defined at a place over
two. Put \(\Gamma=[G_L,G_L]\), with closure understood.
Then \(\Gamma\subset\Delta_0\), and
\[
                       H=P_\ell^{\bar s_P(\Gamma)}
\]
is a closed subgroup preserved by \(G_L\).
For \(\ell\ne2\), choose a local tame inertia element with
nonzero \(\ell\)-adic tame parameter. The actual commutator
of that element with a local Frobenius belongs to \(\Gamma\).
Its first derivation is a nonzero multiple of \(N\):
the multiplier is the tame parameter times \(q-1\), where
\(q>1\) is the residue cardinality. This is the tame
Frobenius relation applied to the based monodromy action.

By (A.35.6), every nonidentity element of \(H\) has first
augmentation degree two, on the line
\(\mathbf Q_\ell\omega_0\). Hence \(H\) is contained in
degree two and meets degree three trivially. Its map to
the degree-two augmentation quotient is an injective
homomorphism with a closed image on a one-dimensional
rational line in a finite free \(\mathbf Z_\ell\)-module.
If nonzero, that image is isomorphic to \(\mathbf Z_\ell\).
Consequently \(H\) is either trivial or procyclic.

Suppose it is nontrivial, with topological generator \(h\).
Its arithmetic conjugation character \(\theta\) occurs
in degree two of the proper curve's homology-generated
filtration. At every good Frobenius it therefore has
weight two. We spell out the incompatible weight-one
detector. Choose a finite \(\ell\)-group quotient \(Q\)
in which \(h\) is nonidentity, and let \(P'\) be the
inverse image of the cyclic subgroup generated by its
image. The homomorphism
\[
              P'\longrightarrow\langle\bar h\rangle
                      \longrightarrow\mathbf Z/\ell
\]
is nonzero on \(h\). Thus \(h\) has nonzero image in
the integral homology of \(P'\), which is torsion-free:
\(P'\) is the pro-\(\ell\) group of a finite proper
\'etale cover, by the induced topology for covers with
\(\ell\)-group Galois closure.

After replacing \(L\) by a finite extension, its arithmetic
section preserves \(P'\). There are only finitely many
subgroups of its fixed index in the finitely generated
pro-\(\ell\) group, so this stabilizer is open. The cover
and this action are then defined over that number field.
The nonzero rational line generated by \(h\) in its
proper homology carries the same character \(\theta\).
It has weight one. At a prime good for both curves, and
different from \(\ell\), the complex absolute values
would have to be both \(q\) and \(q^{1/2}\), a contradiction.
Therefore \(H=1\), proving (A.35.2).

The punctured-to-proper map is equivariant, and its kernel
is the closed normal subgroup generated by the filled
cusp inertias. This proves the final assertion and also
identifies why this proper-image argument alone does not settle
the full punctured fixed-subgroup problem; A.36 supplies that step.
\end{proof}
